% Generated draft: requires mathematical and source review.
\documentclass[11pt]{article}
\usepackage[margin=25mm]{geometry}
\usepackage{fontspec}
\setmainfont{Latin Modern Roman}
\usepackage{amsmath,amssymb,mathtools,unicode-math}
\setmathfont{Latin Modern Math}
\usepackage{enumitem,needspace,xurl,hyperref}
\setlength{\parindent}{0pt}
\setlength{\parskip}{5pt}
\setlength{\emergencystretch}{3em}
\newcommand{\N}{\mathbb{N}}
\newcommand{\Z}{\mathbb{Z}}
\newcommand{\Q}{\mathbb{Q}}
\newcommand{\R}{\mathbb{R}}
\newcommand{\C}{\mathbb{C}}
\newcommand{\F}{\mathbb{F}}
\newcommand{\A}{\mathbb{A}}
\newcommand{\PP}{\mathbb{P}}
\newcommand{\E}{\mathbb{E}}
\newcommand{\eps}{\varepsilon}
\newcommand{\dd}{\,\mathrm{d}}
\begin{document}
\section*{UnsolvedMath ID: 801121}
\textbf{Source problem code:} catalogue entry
\subsection{Definitions and mathematical statement}
Let \(\mathcal R\) denote the separable hyperfinite factor of type \(\mathrm{II}_1\). A factor of type \(\mathrm{II}_1\) is an infinite-dimensional von Neumann algebra with center \(\C 1\) and a faithful normal tracial state \(\tau\) satisfying \(\tau(1)=1\). Such a factor \(P\) is called \emph{McDuff} if
\[
P\cong P\,\overline{\otimes}\,\mathcal R,
\]
where \(\overline{\otimes}\) denotes the spatial tensor product of von Neumann algebras and \(\cong\) denotes von Neumann algebra isomorphism.

Under the type \(\mathrm{II}_1\)-factor interpretation suggested by the source title, the question is: for factors \(M,N_1,N_2\) of type \(\mathrm{II}_1\), does
\[
\left(M\cong N_1\,\overline{\otimes}\,N_2
\quad\text{and}\quad M\text{ is McDuff}\right)
\quad\Longrightarrow\quad
\left(N_1\text{ is McDuff}\quad\text{or}\quad N_2\text{ is McDuff}\right)
\]
always hold?

The supplied excerpt does not explicitly state the type assumptions on the tensor factors or whether separable preduals are required. The displayed formulation therefore records the suggested type \(\mathrm{II}_1\)-factor interpretation; its intended scope, including separability hypotheses, requires source clarification. The introductory assertion that Jolissaint proved certain group von Neumann algebras McDuff refers to an unidentified family, so it cannot supply an additional precise hypothesis here.
\subsection{Short English statement}
If a McDuff factor of type \(\mathrm{II}_1\) decomposes as a tensor product of two factors, must at least one tensor factor be McDuff? The intended domain and any separability assumptions need clarification. The local record labels this question open; that label is supplied metadata, not an independently verified assessment.
\subsection{Sources}
\begin{itemize}
\item UnsolvedMath local registry, permanent ID \(801121\), source problem code \texttt{catalogue entry}; supplied complete source record.
\item Elayavalli, \emph{50 Open Problems: Ultraproduct II\_1 Factors} (2025), Problem 22, as identified in the supplied provenance: \url{https://arxiv.org/abs/2511.20377}. The record attributes the problem to Elayavalli, with proposed year 2025. The supplied provenance records access on 2026-10-16 and extraction by \texttt{arxiv\_latex}; these are supplied provenance claims, not independent source verification.
\item The excerpt cites the label \texttt{[Jolissaint1998]} for the assertion that certain group von Neumann algebras are McDuff. No full bibliographic reference or identification of those groups is supplied.
\end{itemize}
\end{document}
